\documentclass[11pt,a4paper]{article}
\usepackage[margin=0.82in]{geometry}
\usepackage{booktabs}
\usepackage{graphicx}
\usepackage{hyperref}
\usepackage{microtype}
\hypersetup{colorlinks=true,linkcolor=black,urlcolor=blue,citecolor=black}

\title{Choosing a Topic Model for CDP Climate Disclosures, 2016--2024}
\author{}
\date{29 September 2026}

\begin{document}
\maketitle

\begin{abstract}
This note compares three unsupervised approaches to classifying CDP narratives
about emissions-reduction initiatives, climate risks, and climate opportunities:
TF--IDF with nonnegative matrix factorization (NMF), multilingual MiniLM
embeddings with $k$-means, and BERTopic. All methods are tested on the same
year-balanced sample; both neural approaches use the same 256-token embeddings.
The recommended main specification is MiniLM with 12 clusters per response
type, with NMF as a lexical sensitivity check. BERTopic runs successfully, but
its substantial outlier fraction makes it less suitable for complete
company-year coverage in this dataset.
\end{abstract}

\section{Research setting and methodological rationale}
The extracted CDP corpus contains 314,298 public response rows from the
questionnaire items on implemented reduction initiatives, climate risks, and
opportunities across 2016--2024. The questionnaire identifiers and answer
formats change across years, so clustering each response type separately is
more defensible than pooling all three. Topic assignments are descriptions of
disclosures, not measurements of actual abatement or causal policy effects.
For initiatives, the clustering text combines the reported category, type,
and the free-text Comment. The Comment was present in the text used for
the fitted models; the enriched output now also exposes it in a separate
column. A company-year reporting-currency selection has been joined from
CC0.4/C0.4 or, in 2024, Q1.2. It is available for 314,215 of the
314,298 source rows, but was not used as a topic-model feature and does
not convert monetary amounts into a common currency.

Sentence-BERT-style encoders produce fixed-length vectors suited to semantic
comparison of texts \cite{reimers2019}. BERTopic combines transformer
embeddings, dimensionality reduction, density clustering, and class-based
TF--IDF topic words \cite{grootendorst2022}. Studies of short survey
narratives find neural approaches useful, but also show that word-coherence
scores can disagree with expert judgments \cite{xu2022,hoyle2021}. We
therefore compare coverage and semantic separation as well as word-level
interpretability, and retain representative responses for human review.

\section{Common sample and evaluation}
For each response type and disclosure year, we draw 200 distinct narratives
of at least 35 characters after removing exact duplicate texts within type.
The resulting sample has 5,400 records: 1,800 per type, with 1,440 training
and 360 held-out records per type. The random seed is 42. The
\href{https://huggingface.co/sentence-transformers/paraphrase-multilingual-MiniLM-L12-v2}{multilingual
MiniLM-L12-v2} encoder produces normalized 384-dimensional embeddings.
We explicitly set \texttt{max\_seq\_length = 256}; the distributed sentence
transformer defaults to 128 tokens, although its underlying BERT model has
512 position embeddings. Thus the 256-token setting is technically supported
but is an override of the model's distributed pooling configuration, not a
separately validated long-document model.

NMF uses 12 components on word-level TF--IDF. MiniLM with $k$-means uses 12
clusters and class-level TF--IDF words to describe each cluster. BERTopic
uses the same precomputed MiniLM embeddings, a five-dimensional cosine UMAP
projection, and HDBSCAN with a minimum cluster size of 25 and
\texttt{min\_samples = 1}. BERTopic chooses its own number of topics and may
label responses as outliers. All three methods use the same training/test
split. BERTopic was fitted, used to transform held-out narratives, and saved;
the comparison is not based on an import-only smoke test.
For the exclusive assignment diagnostics, an NMF response is assigned to its
highest-weight component, even though NMF itself permits topic mixtures.

We report cosine silhouette on held-out embeddings (higher means more
separation), mean pairwise normalized pointwise mutual information (NPMI) for
the ten leading words per topic (higher means more within-topic word
co-occurrence), and top-ten word diversity (the fraction of distinct words
across topics). NPMI uses the common sampled corpus as its reference, with
zero-co-occurrence pairs set to $-1$. For BERTopic, silhouette is calculated
\emph{only among assigned responses}; this conditional score cannot be read
without its outlier rate. None of these diagnostics supplies ground-truth
topic labels \cite{hoyle2021,hoyle2022}. Silhouette in MiniLM space also
structurally favors the two models that cluster MiniLM embeddings; we
therefore include a parallel TF--IDF-space silhouette check below.

\section{Benchmark results}
\begin{table}[htbp]
\centering
\caption{Held-out topic-model diagnostics by CDP response type}
\label{tab:models}
\small
\setlength{\tabcolsep}{4.0pt}
\begin{tabular}{llrrrrrr}
\toprule
Type & Method & Topics & Outliers (\%) & Sem.\ sil. & Lex.\ sil. & NPMI & Diversity \\
\midrule
Initiatives & TF--IDF + NMF       & 12 & 0.0  &  0.041 & 0.054 & 0.188 & 0.700 \\
            & MiniLM + $k$-means & 12 & 0.0  &  0.124 & 0.048 & 0.220 & 0.608 \\
            & BERTopic           & 18 & 28.3 &  0.134 & 0.067 & 0.096 & 0.667 \\
\addlinespace
Risks       & TF--IDF + NMF       & 12 & 0.0  & -0.020 & 0.032 & 0.296 & 0.792 \\
            & MiniLM + $k$-means & 12 & 0.0  &  0.057 & 0.007 & 0.269 & 0.500 \\
            & BERTopic           & 13 & 26.4 &  0.071 & 0.013 & 0.225 & 0.638 \\
\addlinespace
Opportunities & TF--IDF + NMF       & 12 & 0.0  & 0.016 & 0.020 & 0.162 & 0.792 \\
              & MiniLM + $k$-means & 12 & 0.0  & 0.063 & 0.006 & 0.179 & 0.458 \\
              & BERTopic           & 18 & 58.1 & 0.058 & 0.002 & 0.118 & 0.544 \\
\bottomrule
\end{tabular}

\vspace{3pt}
\begin{minipage}{0.98\linewidth}\footnotesize
\emph{Notes.} Each type has 1,440 training and 360 held-out narratives.
\emph{Sem.} and \emph{Lex.} are cosine silhouette in the common MiniLM and
TF--IDF spaces, respectively. NPMI and word diversity describe topic words,
not assignment accuracy. Both BERTopic silhouettes exclude held-out outliers;
NMF and $k$-means assign all held-out responses.
\end{minipage}
\end{table}

MiniLM with $k$-means has higher held-out semantic separation than NMF for
all three response types and higher NPMI for initiatives and opportunities.
NMF has higher lexical silhouette for all three types, yields more distinct
top words, and has the highest risk-word coherence. These diagnostics expose
an evaluation-space trade-off, not a decisive ordering of topic validity.
BERTopic's conditional silhouette is sometimes high, but 28.3\% of
initiatives, 26.4\% of risks, and 58.1\% of opportunities are unassigned.
Across alternative HDBSCAN settings, a more conservative specification
(minimum cluster size 25, \texttt{min\_samples = 5}) collapses risks into
two topics, while a smaller cluster size (10, \texttt{min\_samples = 2})
creates 30--42 topics per type and leaves roughly 40--64\% unassigned.
The selected BERTopic setting in Table~\ref{tab:models} is the most usable
of these three tested variants, but its coverage remains weak.
On this machine, NMF fitting took 0.2--0.7 seconds per type, $k$-means
0.4--0.5 seconds, and BERTopic 3.2--12.9 seconds. These times exclude the
shared 303.6-second MiniLM embedding pass and are hardware-dependent.

The inspected leading terms are substantively recognizable: initiative
clusters include lighting, HVAC, solar power, waste and materials, and fleet
changes; risk clusters include carbon pricing, weather events, materials
costs, and reputation. Opportunity labels overlap more, particularly around
new products and services. These are observations about a sampled vocabulary,
not a completed expert coding exercise. CDP's selected answer categories are
included in the clustering narrative, so some topics may reproduce the
questionnaire's taxonomy rather than discover entirely new themes.

\section{Recommendation for the decarbonization panel}
Use 12-type-specific MiniLM/$k$-means clusters as the main descriptive theme
variables. They cover all tested responses and avoid BERTopic's variable
outlier selection, which would otherwise make year-to-year company panels
harder to compare. Retain TF--IDF/NMF as a lexical robustness analysis,
especially for risks, and use BERTopic for exploratory discovery among its
well-supported non-outlier groups, not as the sole assignment system.
This is a recommendation for \emph{topic organization}, not a claim that
MiniLM has established causal effects on emissions.

For the full corpus, a computationally cheaper implementation embeds a
year-balanced training sample with MiniLM, fits semantic clusters, and uses
a separately validated TF--IDF classifier to assign the remaining texts.
Those inferred labels must be marked separately from direct embedding labels;
agreement with a held-out MiniLM assignment is a surrogate-fidelity measure,
not evidence of substantive topic validity. Before using cluster prevalence
as a regressor, manually code a stratified sample of response texts from
each topic and year, record inter-rater agreement, and audit the unmatched
or ambiguous narratives. Aggregate labels to the company-year level only
after that check. Reported initiative savings must remain distinct from
observed Scope 1--3 emissions changes.

The 256-token full-corpus hybrid run produced assignments for 312,907 of
314,298 records. Of these, 18,000 were directly embedded and 294,907 were
assigned by the TF--IDF surrogate; 1,391 very short narratives were left
unassigned. On exact-text-grouped holdouts, surrogate accuracy/macro-F1 were
0.772/0.752 for initiatives, 0.687/0.624 for risks, and 0.668/0.671 for
opportunities. The weaker risk and opportunity fidelity warrants additional
manual review and caution in any downstream regression using those labels.
Neither this fidelity test nor the unsupervised benchmark establishes
year-to-year measurement invariance. The full-corpus implementation uses
MiniBatchKMeans and surrogate assignment, so it is a scalable approximation
to, not the identical fitted model from, the common-sample benchmark.
The benchmark manifest, metric table, topic words, representative narratives,
and fitted BERTopic models are saved with the analysis output. Computation
used locally cached model weights without transmitting CDP text to an
external language-model service.

\begin{thebibliography}{9}
\bibitem{reimers2019} Reimers, N. and Gurevych, I. (2019).
Sentence-BERT: Sentence embeddings using Siamese BERT-networks.
\emph{Proceedings of EMNLP-IJCNLP}.
\url{https://aclanthology.org/D19-1410/}.

\bibitem{grootendorst2022} Grootendorst, M. (2022).
BERTopic: Neural topic modeling with a class-based TF-IDF procedure.
\emph{arXiv:2203.05794}.
\url{https://arxiv.org/abs/2203.05794}.

\bibitem{xu2022} Xu, X., Stulp, G., Van Den Bosch, A., and Gauthier, A. (2022).
Understanding narratives from demographic survey data: A comparative study
with multiple neural topic models. \emph{Proceedings of NLP+CSS}, 33--38.
\url{https://aclanthology.org/2022.nlpcss-1.4/}.

\bibitem{hoyle2021} Hoyle, A., Goel, P., Hian-Cheong, A., Peskov, D.,
Boyd-Graber, J., and Resnik, P. (2021).
Is automated topic model evaluation broken? The incoherence of coherence.
\emph{NeurIPS 34}.
\url{https://proceedings.neurips.cc/paper/2021/hash/0f83556a305d789b1d71815e8ea4f4b0-Abstract.html}.

\bibitem{hoyle2022} Hoyle, A., Goel, P., Sarkar, R., and Resnik, P. (2022).
Are neural topic models broken? \emph{Findings of EMNLP}.
\url{https://aclanthology.org/2022.findings-emnlp.390/}.
\end{thebibliography}
\end{document}
